% Copyright (c) 2026 Eugene Joseph M. Ragasa.
% SPDX-License-Identifier: Apache-2.0
% Author: Eugene Joseph M. Ragasa
% Formal proof completed: 2026-10-31
\documentclass[10pt]{article}
\usepackage[T1]{fontenc}
\usepackage{amsmath,amssymb,booktabs,tabularx}
\usepackage[margin=0.75in]{geometry}
\usepackage[hidelinks]{hyperref}
\newcolumntype{Y}{>{\raggedright\arraybackslash}X}
\title{Proof Dossier: Normalized One-Circle Bessel Identity}
\author{Eugene Joseph M. Ragasa}
\date{Formal proof completed: 2026-10-31}
\begin{document}
\raggedright
\maketitle

\section{Identity and status}
\begin{tabularx}{\textwidth}{@{}lY@{}}
\toprule
Statement ID & \texttt{BG-FOURIER-THM-129} \\
Status & proved; no admissions \\
Lean declaration & \path{ProjectKoios.BandGauge.Fourier.integral_exp_oneCircleSinePhase_eq_besselJZero} \\
Owning module & \path{ProjectKoios/BandGauge/Fourier/NormalizedCircleBessel.lean} \\
Author and license & Eugene Joseph M. Ragasa; \texttt{Apache-2.0} \\
\bottomrule
\end{tabularx}

\section{Exact mathematical statement}
Let $\mathbb T_1=\mathbb R/\mathbb Z$ carry its standard Haar probability measure, and define
$s(q)$ through the standard complex-circle map, so that $s([x])=\sin(2\pi x)$. For every real
amplitude $b$,
\[
  \int_{\mathbb T_1}\exp\!\bigl(i b\,s(q)\bigr)\,dq=J_0(b),
\]
where $J_0$ is Mathlib's \texttt{Complex.besselJ 0}. The phase sign is positive and no extra
factor of $2\pi$ is included in $b$.

\section{Source relation and immutable identity}
The full-angle input is \texttt{BG-SPECIAL-THM-014}, whose proof is attributed to the immutable
Apache-2.0 source recorded in its dossier and \path{docs/source-audit.md}. This theorem is a
Project Koios-owned normalization bridge using pinned Mathlib's \texttt{UnitAddCircle} APIs.

The normalization was selected for source-neutral reuse by KSD's 2-D and 3-D flat-parent
comparators. Immutable KSD source identity is commit
\texttt{fcfc4d5fa76b08c040a02c97a109eb5f5fc06e48}, tree
\texttt{fe8b47b33fcf0b24c5b28f8799ed220e3f6d7d20}; the relevant protocol locators and hashes are
recorded in \path{docs/source-audit.md}. This theorem does not assert that KSD's notation or runtime
phase equals the Lean input.

The local module SHA-256 at dossier preparation is
\texttt{0f1a1e808d87302cfa51dbad27dc971fd349a8a9dd40ac063653be46a3f72961}. The Formal base is
merge commit \texttt{7188b05149c598d2edc7ec8eb2af331c54b313ad}; no topic commit is claimed before
checkpoint authorization.

\section{Assumptions and conventions}
The carrier is Mathlib's \texttt{UnitAddCircle}; its standard Haar measure has total mass one. A
real representative $x$ maps to $\sin(2\pi x)$. The exact complex phase is $i b\sin(2\pi x)$.
These choices fix endpoint identification, scaling, normalization, and sign internally, but they do
not authenticate any external convention.

\section{Proof structure}
Use \texttt{UnitAddCircle.intervalIntegral\_preimage} to rewrite the normalized circle average as
an integral over $x\in[0,1]$. Substitute $\theta=2\pi x$, producing a factor $(2\pi)^{-1}$. Shift
$\sin\theta$ to $\cos(\theta-\pi/2)$ and use $2\pi$ periodicity to restore the interval
$[0,2\pi]$. Apply \texttt{BG-SPECIAL-THM-014}; the normalization factor cancels $2\pi$.

\section{Lean realization and trusted base}
The proof uses the standard quotient-to-circle map, \texttt{UnitAddCircle.intervalIntegral\_preimage},
interval scaling and translation, \texttt{Real.cos\_sub\_pi\_div\_two}, and periodic interval
translation. Toolchain: Lean \texttt{4.35.0} at
\texttt{eb6cafb7290e8bd5f2937522d8d804f58f02bafc}; Mathlib
\texttt{56a005855ea17e0b0958f73ec801ae140b0a1679}.
\begin{verbatim}
#print axioms ProjectKoios.BandGauge.Fourier.
  integral_exp_oneCircleSinePhase_eq_besselJZero
-- [propext, Classical.choice, Quot.sound]
\end{verbatim}
No repository-owned axiom or admission occurs.

\section{Downstream exact consequences}
The existing normalized product-torus factorization gives an arbitrary finite cardinality power,
$J_0(a)^2$ on \texttt{Fin 2}, $J_0(b)^3$ on \texttt{Fin 3}, and the exact substitution
$b=2a/3$. Those corollaries remain source-neutral mathematical identities, not authenticated KSD
comparators.

\section{Evidence and interpretation boundaries}
The theorem does not bind an external half-open coordinate, phase formula, amplitude, or measure;
authenticate quadrature, finite meshes, runtime values, binary64, or serialization; establish
finite-to-continuum convergence; or certify a scientific comparator, model, certificate, or
publication claim. The checked runtime-certificate count remains zero.

\section{AI, automation, and review record}
The Lean bridge and dossier were AI-assisted and checked by the maintainer. Kernel checking is not
independent mathematical, source-fidelity, numerical, scientific, novelty, or publication review.

\section{Review and revision history}
\begin{tabularx}{\textwidth}{@{}lllY@{}}
\toprule Date & Revision & Process & Outcome and limitation \\
\midrule 2026-10-31 & initial dossier & AI-assisted maintainer review & normalized identity proved; external/runtime binding and independent patch review remain open \\
\bottomrule
\end{tabularx}
\end{document}
